\section{Discussion}

In this HMD reading dataset, the short-term repeatability of an exploratory interocular reading-event contrast depended strongly on how its events were constructed. On a common cohort of 376 participants (floor 5, R1 S1--S2), keeping only the mean-speed events that a matched per-eye detection overlapped raised ICC(A,1) from .374 to .693. The windows, peaks and classes of the kept events were unchanged, so the increase came from removing events, not from measuring the rest differently. Most of it came from one participant: without that participant, the unfiltered ICC was .639, the admitted ICC .693, and the contrast .054 [$-$.007, .129]. The admitted estimate was thus insensitive to that participant, whereas the unfiltered estimate depended on it. A conventional 1,000$^\circ$/s peak-velocity ceiling, applied post hoc, gave a similar point estimate (.711 against .693; admitted minus ceiling $-$.018 [$-$.090, .039]), and the low unfiltered estimate itself depended on detector settings, ranging from .226 to .703 across velocity thresholds, smoothing windows and pairing gaps, whereas point estimates under every plausibility rule stayed between .66 and .74, each variant in its own eligible cohort. In the units of $A$, the median absolute session difference was about .02 under all constructions (a multiplicative change of about 4\% in $e^{2A}$) and about 11.5\% lower with admission than without. After admission, replacing admitted events with matched per-eye events, or adding recording-centred abstention, changed the estimate by smaller amounts with uncertain increments. At floor 5 for R1 S1--S2, which recordings were eligible moved the mean\_u ICC(A,1) from .170 in its own cohort ($n$ = 379) to .375 in the unfiltered/centred intersection ($n$ = 377), a difference of two participants, one of them with one extreme session. At floor 10, the participant who carried the admission contrast was outside the joint cohort. Admission, eligibility and the floor thus acted largely through the same few extreme sessions. The unfiltered estimate was fragile in a specific way. A few participants accounted for most of the squared session differences, the Pearson correlation of unfiltered with other constructions was moderate while their rank correlation was high, and the eligibility floor changed the size of the increments, at least partly by including or excluding a few recordings with extreme $A$. The absolute-agreement ICC is sensitive to a few extreme recording values, and in this dataset large recording-level discrepancies and leverage on the ICC were concentrated in a few participants under the unfiltered mean-speed construction.

For reliability reports of binocular kinematic markers, these results suggest stating, alongside the ICC, the detector, the signal it thresholds, its threshold, minimum run length and velocity filter; the binocular pairing rule; the rule for admitting events and any plausibility ceiling; any abstention rule and the pool over which its reference is computed; the exclusion windows and class definitions; the eligibility floor, with the ICC at one alternative floor; the cohort on which compared constructions are evaluated; the ICC form, unit and interval; and sample-level data-quality metrics. An influence analysis (the largest single-participant share and the leave-one-out range of the ICC) and a rank correlation make it visible when a few participants carry the estimate. For this contrast and pipeline with the frozen detector settings, admitting only per-eye-supported events gave higher short-term repeatability than admitting all mean-speed events, chiefly by removing the events behind one participant's extreme session difference. Many of those events were physiologically implausible: they combined a large left-eye displacement with little right-eye displacement and left-eye peak speeds above the 1,000$^\circ$/s commonly used as a plausibility ceiling~\citep{nystrom2010}. The mean-speed detector applied no such ceiling and accepts an event driven by one eye alone, so it is permissive by design. Admission corresponds to requiring that both eyes register the movement, a binocular criterion analogous to the temporal-overlap requirement used to define binocular microsaccades~\citep{engbert2003}. A peak-velocity ceiling or a binocular requirement would therefore be a natural part of a baseline construction. The post hoc ceiling comparator supports this reading. The ceiling and admission removed largely different events, yet the ceiling alone gave a similar ICC, adding admission to it changed the ICC by .005, and admitted minus ceiling stayed between $-$.062 and .014 across ceilings of 750 to 1,500$^\circ$/s and across detector settings, with the only interval excluding zero favouring the ceiling. The admission contrast reported here is therefore the cost of omitting a plausibility safeguard from a permissive baseline, not a gain specific to admission; a ceiling is the simpler safeguard, whereas admission additionally requires a per-eye detector. Admission also does not guarantee robustness: in the R1--R3 unfiltered own cohorts, an extreme session whose return sweeps all had per-eye support remained after admission (Results). Admission without centring yielded ICC(A,1) = .693, and adding centring changed it by .034 [$-$.008, .100]. Centring, admission and per-eye detection reached similar ICCs on the joint cohort, and their recording-level estimates agreed closely where centring was applied. Apart from the implausible events of that session, higher repeatability does not establish that the admitted events are more valid, and the rejected events in general were not established to be artefacts. The per-eye detector reads the same signals as the mean-speed detector and is not independent ground truth. Its agreement with centring on which uncontaminated events to reject (odds ratio 110.8; Results) shows that the two rules flag overlapping events, not that those events are errors, and it is partly mechanical: an event driven by one eye tends both to have an extreme ratio and to lack a matched per-eye pair. Events without per-eye support could include genuine but atypical saccades, such as saccades with large transient interocular differences, alongside tracking errors.

The synthetic injections tested centring applied to mean-speed events; they did not test admission or per-eye detection, so they do not show which faults admission removes. Within the limits of the injected fault models, they clarify what centring can detect. One-eye step-like displacements of 4$^\circ$ or 8$^\circ$ during fixation produced extreme ratios and were almost always rejected once they reached the centring step. Among injections reaching centring, recall was .772 for one-eye slowing by a factor of 3 and .376 for brief 6$^\circ$ peak glitches, below the frozen .80 criterion. For factor-3 slowing, 50.4\% of injections reached centring and 46.7\% were excluded by the duration rule; these earlier exclusions are separate pipeline outcomes, not centring misses. Mild within-event displacement scaling increased rejection of the same target events by at most .009; because it changes $l$ by only about .15, well inside the $\log 2$ band, this specificity control is also close to a sanity check. The rule therefore acts on gross one-eye ratio departures and is not a general detector of one-eye faults. In a mixed synthetic contamination, the full-sample ICC changed little, but the estimate for the exposed participants fell under both the centred and unfiltered constructions. Only 17 of the 376 evaluated participants were exposed and evaluable, so a full-sample ICC can change little even when the exposed recordings are affected. The large one-eye steps that met the criterion act mainly as a sanity check, since a step of several degrees in one eye produces a ratio far outside the $\log 2$ band. We therefore do not recommend centring as a fault filter; it is better treated as a descriptive indicator of events whose ratio departs from the recording's own median.

The positive sign of the amplitude-matched left/right difference $D$ in both reading and random-saccade tasks, and its association with $A$ across sessions, agree with the classic observation that the abducting eye moves faster in horizontal saccades~\citep{collewijn1988}. They do not identify the source of $A$. In $A$, leftward return sweeps are larger and serve a different role in reading than rightward progressive saccades, and the matched controls equate amplitude but not task role or, for TEX, the full range of amplitudes. Left--right asymmetries in saccade latency, gain and peak velocity have also been examined in their own right~\citep{vergilino2012}, and direction-specific influences of this kind were not modelled here. Other work has assessed interocular saccade-velocity measures under different conditions. The DEMoNS protocol standardizes infrared-oculography saccade measurements and reports abducting/adducting peak-velocity ratios~\citep{nijbijvank2018}. Individual differences in the asymmetry of binocular saccades have also been analysed with mixed-effects models~\citep{kloke2006}. Those studies differ from ours in estimator (neither uses a median log-ratio contrast between reading-event classes), in population and recording device, in the interval between measurements, and in aggregation. Similar ICC values across them would therefore not imply comparable measurement performance. In the same dataset, binocular features have been assessed for eye-movement biometrics~\citep{yao2026}. That work asks which features identify individuals; ours asks how event construction changes the reliability of one interocular marker, so the two are complementary. Saccade main-sequence parameters and other oculomotor measures have been reported to be repeatable within individuals~\citep{gibaldi2021,bargary2017}. By common guidelines, the short-term ICC point estimates of about .69 to .73 obtained with admitted, per-eye or centred events indicate moderate reliability, with 95\% CIs extending into the good range, whereas the joint-cohort estimate for unfiltered mean-speed events (.374) falls in the poor range~\citep{koo2016}.

Twelve limitations bound these conclusions. First, the analysis is exploratory and post hoc; the comparison of unfiltered with centred mean-speed events was part of the frozen specification, but its joint-cohort evaluation and the admission contrast were added afterwards. A pilot on the same reading task informed the estimator, the abstention rule and the eligibility floor, and the event-construction framing was adopted after the results were seen. The peak-velocity ceiling, the absolute session differences and the detector-setting variants were added after review. The contrasts reported here are descriptive of this dataset and need confirmatory replication; in particular, the admission contrast depended largely on one participant and on the detector settings. Second, there is no independent ground truth for which events are valid. We did not use human annotation, and the per-eye detector reads the same signals as the mean-speed detector, so admission, centring and their agreement cannot establish event validity. Third, the abstention rule did not meet its frozen injection criteria for one-eye slowing and peak glitches, so it cannot be relied on to remove these faults, and the injections did not evaluate admission. Fourth, the longitudinal cohorts were small and selected (29 to 56 participants who returned for later rounds), so the longitudinal estimates are imprecise and may not represent the full R1 sample. Fifth, random-saccade $D$ was associated with reading $A$, but its repeatability could not be estimated because no participant had enough pairs in each session. Sixth, direction, amplitude range and event class co-vary within $A$, so its physiological source is not identified. Seventh, all data come from one dataset of college-aged participants, recorded with one HMD eye tracker at 250\,Hz with a chin rest; the two R1 sessions used different text excerpts, and calibration was validated over $\pm 15^\circ$ horizontally~\citep{lohr2023}, narrower than the reading region, so whether the dependence on event construction holds for other devices, sampling rates and filters remains to be tested. Post-saccadic oscillation size varies with saccade size~\citep{hooge2015}, and threshold-based event detection can be sensitive to noise~\citep{ivtadapt2025}. Detector settings were varied one at a time over a narrow range around the frozen values, and other velocity filters and detector families were not tested. Eighth, robustness to calibration was tested only for constant affine per-eye gains. Position-dependent or time-varying calibration errors were not tested. Ninth, repeatability is not validity. A marker can be repeatable because it captures a stable property of the measurement pipeline or device fit, not of the oculomotor system. Tenth, $A$ is not a pure velocity measure. At low sampling rates and with low-pass filtering, peak velocity is attenuated, more strongly for small saccades~\citep{mack2017}, so a log ratio of peak velocities can carry interocular differences in amplitude, filtering and viewing geometry, and the return-sweep and progressive classes differ in amplitude. Apart from the matched $D$, we did not control $A$ for amplitude, and post-saccadic oscillations were not detected separately. Eleventh, centring bounds $A$ by construction. Every kept event lies within $\log 2$ of the recording median, so $|A|$ is below $\log 2$ (about .69). Part of the higher ICC of centred constructions may reflect this truncation of extreme values, although admission without centring reached a similar ICC; in samples with genuinely large, class-specific interocular asymmetry the same bound could censor the signal of interest. Twelfth, ICC(A,1) is not robust to extreme values, and apart from absolute session differences of $A$ we did not report a robust ICC, the standard error of measurement or limits of agreement adapted to repeated measures; the pooled limits of agreement between constructions weight participants by their number of recordings, and the many contrasts reported carry unadjusted intervals.

Two steps would test these findings. A confirmatory replication should freeze the constructions, with a conventional peak-velocity ceiling as the baseline, together with the floors, detector settings and contrasts reported here, and apply them to binocular HMD reading data from another device. An annotated reference set of binocular saccades, including events without per-eye support, would show what the rejected events are. Reliability reports for interocular saccade markers in HMD data should state the event-construction and eligibility rules that produced their estimates.
